(1) 两个无相互作用的自旋 \(1/2\) 全同费米子被限制在长度 \(L\) 的一维无限深方势阱，单粒子空间能级 \(\varepsilon_n=\pi^2\hbar^2n^2/(2mL^2)\)。外磁场 \(B\ge0\) 导致Zeeman项 \(-\mu B\sigma_z\)，\(\mu>0\)。求前几个单粒子自旋能级与允许的双粒子填充。

(2) 求零温基态在低场和高场的能量，确定从两个反向自旋占同一空间轨道变为同向自旋占两个轨道的临界磁场。

(3) 求基态磁化强度 \(M=-\partial_BE_G\)，说明跳变的物理意义，并讨论临界点的简并。

(4) 计算两个区间内的力学壁压力 \(P=-\partial_LE_G\)，比较自旋极化前后的压强，讨论 \(L\) 改变时阈值的标度。
